% GENERATED FILE. Edit canonical lessons, then run npm run generate:books.
% canonical-source-hash: fnv1a64:7ad0fc9a38579274
% canonical-lessons: PA-C341-shakk, PA-C341-yakin, PA-C341-bharosa, PA-C341-andaza, PA-C341-jachna

\chapter{I'm Not So Sure}
\label{ch:pa-a2-i-m-not-so-sure}
\hlchaptermodality{341}

\begin{chapteropening}
\textbf{By the end of this chapter:} \emph{I can say that I doubt something or that I am sure of it.}

Complete the last lesson of chapter 341: I can say that I doubt something or that I am sure of it.
\end{chapteropening}

\section[\texorpdfstring{\pa{ਸ਼ੱਕ} (shakk)}{ਸ਼ੱਕ (shakk)}]{\pa{ਸ਼ੱਕ} (shakk) — doubt; suspicion}

\label{lesson:PA-C341-shakk}



\begin{quote}
Before the new one: say the Punjabi for to admit, then the Punjabi for a compromise.
\end{quote}

\subsection*{You'll want to know: \pa{ਸ਼ੱਕ}}
\textbf{\pa{ਸ਼ੱਕ}} — \emph{shakk} — \textquotedblleft{}doubt; suspicion\textquotedblright{}.

\begin{quote}
\pa{ਮੈਨੂੰ} \pa{ਤੁਹਾਡੀ} \pa{ਗੱਲ} \pa{ਉੱਤੇ} \pa{ਕੋਈ} \pa{ਸ਼ੱਕ} \pa{ਨਹੀਂ} \pa{ਹੈ।}
\end{quote}

\begin{quote}
\emph{mainū̃ tuhāḍī gall utte koī shakk nahīṁ hai.}
\end{quote}

\begin{quote}
I have no doubt about what you say.
\end{quote}

A masculine noun.

\begin{cousinweb}
\textbf{\pa{ਸ਼ੱਕ}} (\emph{shakk}) is a loan, usually traced through Persian to Arabic \emph{shakk}, \textquotedblleft{}doubt\textquotedblright{}.
\end{cousinweb}

\begin{grammarlens}[title={What you've built}]
One more word for asking and answering.
\end{grammarlens}

\subsection*{Your turn}
\begin{itemize}
\raggedright
  \item \emph{Say it:} \emph{shakk}
  \item \emph{Say it:} \emph{shakk}, once more
  \item \emph{Say it:} \emph{mainū̃ tuhāḍī gall utte koī shakk nahīṁ hai.} — I have no doubt about what you say.
  \item \emph{Recall:} say the Punjabi for to admit, then the Punjabi for a compromise, then say \emph{shakk} again
\end{itemize}

\begin{tcolorbox}[breakable,colback=teal!4,colframe=teal!35!black,title={Before you move on}]
What does \pa{ਸ਼ੱਕ} mean? (\textquotedblleft{}Doubt; suspicion\textquotedblright{}.) Say it once more.
\end{tcolorbox}
\section[\texorpdfstring{\pa{ਯਕੀਨ} (yakīn)}{ਯਕੀਨ (yakīn)}]{\pa{ਯਕੀਨ} (yakīn) — certainty; belief}

\label{lesson:PA-C341-yakin}



\begin{quote}
Before the new one: say the Punjabi for a compromise, then the Punjabi for doubt.
\end{quote}

\subsection*{You'll want to know: \pa{ਯਕੀਨ}}
\textbf{\pa{ਯਕੀਨ}} — \emph{yakīn} — \textquotedblleft{}certainty; belief\textquotedblright{}.

\begin{quote}
\pa{ਮੈਨੂੰ} \pa{ਉਸ} \pa{ਦੀ} \pa{ਗੱਲ} \pa{ਉੱਤੇ} \pa{ਯਕੀਨ} \pa{ਨਹੀਂ} \pa{ਹੈ।}
\end{quote}

\begin{quote}
\emph{mainū̃ us dī gall utte yakīn nahīṁ hai.}
\end{quote}

\begin{quote}
I don't believe what he says.
\end{quote}

A masculine noun.

\begin{cousinweb}
\textbf{\pa{ਯਕੀਨ}} (\emph{yakīn}) is a loan, usually traced through Persian to Arabic \emph{yaqīn}, \textquotedblleft{}certainty\textquotedblright{}.
\end{cousinweb}

\begin{grammarlens}[title={What you've built}]
One more word for asking and answering.
\end{grammarlens}

\subsection*{Your turn}
\begin{itemize}
\raggedright
  \item \emph{Say it:} \emph{yakīn}
  \item \emph{Say it:} \emph{yakīn}, once more
  \item \emph{Say it:} \emph{mainū̃ us dī gall utte yakīn nahīṁ hai.} — I don't believe what he says.
  \item \emph{Recall:} say the Punjabi for a compromise, then the Punjabi for doubt, then say \emph{yakīn} again
  \item \emph{Write it:} copy \textbf{\pa{ਯਕੀਨ}} once, checking the whole word
\end{itemize}

\begin{tcolorbox}[breakable,colback=teal!4,colframe=teal!35!black,title={Before you move on}]
What does \pa{ਯਕੀਨ} mean? (\textquotedblleft{}Certainty; belief\textquotedblright{}.) Say it once more.
\end{tcolorbox}
\section[\texorpdfstring{\pa{ਭਰੋਸਾ} (bharosā)}{ਭਰੋਸਾ (bharosā)}]{\pa{ਭਰੋਸਾ} (bharosā) — trust; reliance}

\label{lesson:PA-C341-bharosa}



\begin{quote}
Before the new one: say the Punjabi for doubt, then the Punjabi for certainty.
\end{quote}

\subsection*{You'll want to know: \pa{ਭਰੋਸਾ}}
\textbf{\pa{ਭਰੋਸਾ}} — \emph{bharosā} — \textquotedblleft{}trust; reliance\textquotedblright{}.

\begin{quote}
\pa{ਕੀ} \pa{ਤੁਸੀਂ} \pa{ਮੇਰੇ} \pa{ਉੱਤੇ} \pa{ਭਰੋਸਾ} \pa{ਕਰਦੇ} \pa{ਹੋ}?
\end{quote}

\begin{quote}
\emph{kī tusī̃ mere utte bharosā karde ho?}
\end{quote}

\begin{quote}
Do you trust me?
\end{quote}

A masculine noun.

\begin{cousinweb}
\textbf{\pa{ਭਰੋਸਾ}} (\emph{bharosā}) is inherited from a reconstructed \emph{*bharavaśya}, \textquotedblleft{}dependence on support\textquotedblright{}; CDIAL lists Punjabi \emph{bharossā}, \textquotedblleft{}confidence\textquotedblright{}, under it (CDIAL 9398).
\end{cousinweb}

\begin{grammarlens}[title={What you've built}]
One more word for asking and answering.
\end{grammarlens}

\subsection*{Your turn}
\begin{itemize}
\raggedright
  \item \emph{Say it:} \emph{bharosā}
  \item \emph{Say it:} \emph{bharosā}, once more
  \item \emph{Say it:} \emph{kī tusī̃ mere utte bharosā karde ho?} — Do you trust me?
  \item \emph{Recall:} say the Punjabi for doubt, then the Punjabi for certainty, then say \emph{bharosā} again
  \item \emph{Write it:} copy \textbf{\pa{ਭਰੋਸਾ}} once, checking the whole word
\end{itemize}

\begin{tcolorbox}[breakable,colback=teal!4,colframe=teal!35!black,title={Before you move on}]
What does \pa{ਭਰੋਸਾ} mean? (\textquotedblleft{}Trust; reliance\textquotedblright{}.) Say it once more.
\end{tcolorbox}
\section[\texorpdfstring{\pa{ਅੰਦਾਜ਼ਾ} (andāzā)}{ਅੰਦਾਜ਼ਾ (andāzā)}]{\pa{ਅੰਦਾਜ਼ਾ} (andāzā) — a guess; an estimate}

\label{lesson:PA-C341-andaza}



\begin{quote}
Before the new one: say the Punjabi for certainty, then the Punjabi for trust.
\end{quote}

\subsection*{You'll want to know: \pa{ਅੰਦਾਜ਼ਾ}}
\textbf{\pa{ਅੰਦਾਜ਼ਾ}} — \emph{andāzā} — \textquotedblleft{}a guess; an estimate\textquotedblright{}.

\begin{quote}
\pa{ਮੇਰਾ} \pa{ਅੰਦਾਜ਼ਾ} \pa{ਗ਼ਲਤ} \pa{ਨਹੀਂ} \pa{ਹੈ।}
\end{quote}

\begin{quote}
\emph{merā andāzā ghalat nahīṁ hai.}
\end{quote}

\begin{quote}
My guess is not wrong.
\end{quote}

A masculine noun.

\begin{cousinweb}
\textbf{\pa{ਅੰਦਾਜ਼ਾ}} (\emph{andāzā}) is a loan, usually traced to Persian \emph{andāza}, \textquotedblleft{}measure\textquotedblright{}.
\end{cousinweb}

\begin{grammarlens}[title={What you've built}]
One more word for asking and answering.
\end{grammarlens}

\subsection*{Your turn}
\begin{itemize}
\raggedright
  \item \emph{Say it:} \emph{andāzā}
  \item \emph{Say it:} \emph{andāzā}, once more
  \item \emph{Say it:} \emph{merā andāzā ghalat nahīṁ hai.} — My guess is not wrong.
  \item \emph{Recall:} say the Punjabi for certainty, then the Punjabi for trust, then say \emph{andāzā} again
  \item \emph{Write it:} copy \textbf{\pa{ਅੰਦਾਜ਼ਾ}} once, checking the whole word
\end{itemize}

\begin{tcolorbox}[breakable,colback=teal!4,colframe=teal!35!black,title={Before you move on}]
What does \pa{ਅੰਦਾਜ਼ਾ} mean? (\textquotedblleft{}A guess; an estimate\textquotedblright{}.) Say it once more.
\end{tcolorbox}
\section[\texorpdfstring{\pa{ਜਚਣਾ} (jachṇā)}{ਜਚਣਾ (jachṇā)}]{\pa{ਜਚਣਾ} (jachṇā) — to seem right; to convince}

\label{lesson:PA-C341-jachna}



\begin{quote}
Before the new one: say the Punjabi for trust, then the Punjabi for a guess.
\end{quote}

\subsection*{You'll want to know: \pa{ਜਚਣਾ}}
\textbf{\pa{ਜਚਣਾ}} — \emph{jachṇā} — \textquotedblleft{}to seem right; to convince\textquotedblright{}.

\begin{quote}
\pa{ਮੈਨੂੰ} \pa{ਉਸ} \pa{ਦੀ} \pa{ਗੱਲ} \pa{ਨਹੀਂ} \pa{ਜਚਦੀ} \pa{ਹੈ।}
\end{quote}

\begin{quote}
\emph{mainū̃ us dī gall nahīṁ jachdī hai.}
\end{quote}

\begin{quote}
What he says doesn't convince me.
\end{quote}

Past (it …): \pa{ਉਹ} \pa{ਜਵਾਬ} \pa{ਮੈਨੂੰ} \pa{ਜਚਿਆ}, \emph{oh javāb mainū̃ jachiā}.

\begin{cousinweb}
\textbf{\pa{ਜਚਣਾ}} (\emph{jachṇā}) is inherited from Sanskrit \emph{yācyate}, \textquotedblleft{}is asked\textquotedblright{}; CDIAL lists Punjabi \emph{jacāuṇā}, \textquotedblleft{}to cause to be examined\textquotedblright{}, under it and notes \emph{jacṇā}, \textquotedblleft{}to be tested\textquotedblright{}, formed from it (CDIAL 10449). CDIAL does not give the everyday sense \textquotedblleft{}to seem right\textquotedblright{}.
\end{cousinweb}

\begin{grammarlens}[title={What you've built}]
One more word for asking and answering.
\end{grammarlens}

\subsection*{Your turn}
\begin{itemize}
\raggedright
  \item \emph{Say it:} \emph{jachṇā}
  \item \emph{Say it:} \emph{jachṇā}, once more
  \item \emph{Say it:} \emph{mainū̃ us dī gall nahīṁ jachdī hai.} — What he says doesn't convince me.
  \item \emph{Recall:} say the Punjabi for trust, then the Punjabi for a guess, then say \emph{jachṇā} again
  \item \emph{Write it:} copy \textbf{\pa{ਜਚਣਾ}} once, checking the whole word
\end{itemize}

\begin{tcolorbox}[breakable,colback=teal!4,colframe=teal!35!black,title={Before you move on}]
What does \pa{ਜਚਣਾ} mean? (\textquotedblleft{}To seem right; to convince\textquotedblright{}.) Say it once more.
\end{tcolorbox}
